\chapter{Conclusion and Future Work}

\section{Conclusion}
Qauth was developed to explore whether a modern authentication platform can be strengthened by combining conventional identity verification with quantum-inspired key generation and behaviour-aware security analysis. The implementation demonstrates that such a combination is both technically feasible and architecturally meaningful. Rather than treating authentication as a one-step password check, Qauth models it as a layered process involving credential validation, secure key lifecycle management, challenge-response proof verification, anomaly interpretation, and administrative visibility.

From a software engineering perspective, the project successfully integrates a React frontend, a Django REST backend, PostgreSQL-based persistence, Redis-backed transient counters, AES-GCM encrypted key storage, JWT-based session handling, and rule-driven anomaly detection into a coherent system. From a security perspective, the project shows that high-entropy user-specific secret material can be incorporated into the authentication workflow in a controlled and observable manner. From an academic perspective, the project provides a useful prototype for understanding how quantum-aware design concepts can influence real application architecture even before full-scale quantum infrastructure becomes commonplace.

\section{Summary of Technical Contributions}
The main technical contributions of Qauth can be summarized as follows:

\begin{enumerate}
    \item development of a custom authentication platform with a dedicated user model and security state tracking,
    \item implementation of a quantum-derived key pipeline using ANU QRNG, BB84 simulation, and secure fallback randomness,
    \item encryption of stored key material using AES-256-GCM and management of reveal and rotation metadata,
    \item implementation of a client-server HMAC challenge-response flow using Web Crypto and backend proof verification,
    \item design of a login anomaly engine using rule-based scoring and Redis-backed sliding-window counters,
    \item creation of audit logging and administrative monitoring endpoints for security observability,
    \item delivery of a full-stack prototype suitable for technical demonstration, extension, and future research.
\end{enumerate}

\section{Achievement of Project Objectives}
The objectives defined at the beginning of the report have largely been achieved. The project successfully implements secure registration, credential login, quantum-key provisioning, challenge-response verification, account locking, behavioural scoring, and administrative monitoring. It also meets the broader objective of producing a technically rich and explainable security system rather than a minimal login page.

However, objective achievement must be interpreted with prototype realism. While the major architectural goals have been met, certain production-level refinements such as extensive automated tests, fully automated alert escalation, and stronger replay-hardening mechanisms remain open for future work. This does not reduce the value of the implementation; rather, it clarifies the stage of maturity of the system.

\section{Practical Significance of Qauth}
Qauth is significant because it demonstrates how emerging security ideas can be translated into a practical application rather than remaining at the level of theory. The project is particularly valuable in educational, research, and prototype contexts where it is important to show:

\begin{enumerate}
    \item how quantum-derived randomness can be integrated into application workflows,
    \item how secret material should be protected throughout its lifecycle,
    \item how login events can be enriched with behavioural context,
    \item how security monitoring can be exposed through administrator-friendly interfaces.
\end{enumerate}

These contributions make the platform useful not only as a software product but also as a teaching and demonstration system for advanced authentication design.

\section{Limitations of the Current Implementation}
The present version of Qauth has several limitations that should be acknowledged clearly:

\begin{enumerate}
    \item the BB84 module is a software simulation and not a hardware QKD implementation,
    \item the quantum challenge mechanism can be strengthened with nonce persistence and replay protection,
    \item automated unit, integration, and performance test suites are limited,
    \item the anomaly engine uses explainable heuristics rather than adaptive machine learning,
    \item threat alert creation can be made more automatic and policy-driven,
    \item production-hardening concerns such as distributed scaling, secret rotation, and hardened deployment remain to be completed.
\end{enumerate}

These limitations are consistent with the project's role as a strong prototype and research-oriented demonstration rather than a final enterprise identity platform.

\section{Future Enhancements}

\subsection{Integration with Hardware Quantum Random Number Sources}
One natural extension is to connect the platform to hardware-based quantum entropy devices or institutionally hosted QRNG services. This would reduce dependence on software simulation and strengthen the authenticity of the quantum-randomness claim.

\subsection{Advanced Post-Quantum Cryptographic Support}
Future versions of Qauth can incorporate standardized post-quantum key encapsulation and digital signature schemes. This would move the platform from quantum-aware design toward stronger end-to-end resistance against future quantum adversaries.

\subsection{Real-Time Threat Intelligence Integration}
The anomaly engine can be expanded by integrating IP reputation feeds, device intelligence, geo-velocity analysis, and adaptive risk thresholds. This would allow the system to evolve from a rule-based prototype into a richer real-time threat evaluation framework.

\subsection{Scalable Cloud Deployment}
The system can be containerized and deployed through cloud-native infrastructure with separate API, database, cache, and worker services. This would support horizontal scaling, centralized logging, and more realistic production simulation.

\subsection{Expanded Analytics and Reporting}
The administrator dashboard can be extended with longitudinal analytics, user-behaviour baselines, alert trend charts, exportable security reports, and forensics-oriented drill-down interfaces. Such additions would improve the platform's operational value.

\section{Final Remarks}
Qauth demonstrates that authentication can be made significantly richer when cryptographic discipline, quantum-inspired key generation, behavioural analysis, and visibility are treated as parts of one system. The project does not claim to solve the entire challenge of future-proof authentication, but it establishes a practical and well-structured foundation for that journey. As quantum computing and post-quantum security continue to shape the cybersecurity landscape, systems like Qauth provide a timely model for how application developers can begin designing with that future in mind.
